\documentclass[11pt,a4paper]{article}

\usepackage[utf8]{inputenc}
\usepackage[indonesian]{babel}
\usepackage[a4paper,top=2cm,bottom=2cm,left=2cm,right=2cm]{geometry}
\usepackage{amsmath,amsfonts,amssymb}
\usepackage{graphicx}
\usepackage{booktabs}
\usepackage{enumitem}
\usepackage{xcolor}
\usepackage{colortbl}
\usepackage{tcolorbox}
\usepackage{fancyhdr}
\usepackage{titlesec}
\usepackage{listings}
\usepackage{tikz}
\usepackage[breaklinks=true,colorlinks=true,linkcolor=primaryblue,citecolor=primaryblue,urlcolor=primaryblue]{hyperref}

% Color Palette (Deep Navy & Accent Gold/Blue)
\definecolor{primaryblue}{RGB}{24, 43, 73}
\definecolor{secondaryblue}{RGB}{41, 128, 185}
\definecolor{accentgold}{RGB}{211, 84, 0}
\definecolor{boxbg}{RGB}{245, 248, 252}
\definecolor{codebg}{RGB}{240, 243, 246}
\definecolor{bordergray}{RGB}{200, 210, 220}
\definecolor{darkgray}{RGB}{44, 62, 80}
\definecolor{codegreen}{RGB}{39, 174, 96}

% Header & Footer
\pagestyle{fancy}
\fancyhf{}
\fancyhead[L]{\small\textbf{\color{primaryblue}Problem Set Minggu 1: Karakteristik \& Dekomposisi Deret Waktu}}
\fancyhead[R]{\small\color{gray}TEE-412 / Teknik Elektro FT UNIB}
\fancyfoot[L]{\footnotesize\color{gray}Dosen: Ir. Novalio Daratha, S.T., M.Sc., Ph.D.}
\fancyfoot[R]{\footnotesize Halaman \thepage\ dari 3}
\renewcommand{\headrulewidth}{0.5pt}
\renewcommand{\footrulewidth}{0.4pt}

% Section format
\titleformat{\section}{\color{primaryblue}\normalfont\large\bfseries}{\thesection.}{0.5em}{}[\titlerule]
\setlist{nosep}

% Listings setup
\lstset{
    language=Python,
    backgroundcolor=\color{codebg},
    basicstyle=\ttfamily\scriptsize\color{darkgray},
    keywordstyle=\bfseries\color{secondaryblue},
    stringstyle=\color{accentgold},
    commentstyle=\itshape\color{codegreen},
    showstringspaces=false,
    breaklines=true,
    frame=single,
    rulecolor=\color{bordergray},
    numbers=none
}

\begin{document}

% HEADER DOKUMEN
\begin{center}
    {\large \textbf{\color{primaryblue}UNIVERSITAS BENGKULU -- FAKULTAS TEKNIK}}\\[0.1cm]
    {\large \textbf{\color{secondaryblue}JURUSAN TEKNIK ELEKTRO}}\\[0.2cm]
    {\LARGE \textbf{\color{primaryblue}PROBLEM SET 1: ANALISIS DERET WAKTU DAN PERAMALAN}}\\[0.1cm]
    {\normalsize \textbf{Materi: Karakteristik Sinyal Temporal, Komponen Dinamis, dan Dekomposisi Klasik}}
\end{center}

\vspace{0.1cm}

% BOX IDENTITAS MAHASISWA & ATURAN
\begin{tcolorbox}[colback=boxbg,colframe=primaryblue,arc=2mm,title=\textbf{Informasi Tugas \& Identitas Mahasiswa}]
\small
\begin{tabular}{p{2.2cm}p{5.5cm}p{2.2cm}p{5.0cm}}
\textbf{Nama} & : \dotfill & \textbf{NPM} & : \dotfill \\
\textbf{Prodi/Kelas} & : S1 Teknik Elektro & \textbf{Hari/Tanggal} & : \dotfill \\
\textbf{Dosen} & : Ir. Novalio Daratha, Ph.D. & \textbf{Batas Waktu} & : 1 Minggu \\
\end{tabular}
\vspace{0.15cm}
\hrule
\vspace{0.15cm}
\textbf{Petunjuk Umum:}
\begin{itemize}[leftmargin=1.2em]
    \item Soal disusun berjenjang berdasarkan \textbf{Taksonomi Bloom (Level C1 sampai C6)} untuk mengukur capaian \textbf{Sub-CPMK 1}.
    \item Tuliskan jawaban Anda secara runtut, jelas, dan sertakan turunan matematis/langkah analitis jika diminta.
\end{itemize}
\end{tcolorbox}

\vspace{0.2cm}

% =============================================================
% SOAL PROBLEM SET BERJENJANG C1 - C6
% =============================================================

\begin{enumerate}[leftmargin=1.5em, label=\textbf{Soal \arabic*.}]

    % ---------------------------------------------------------
    % C1: MENGINGAT (REMEMBERING)
    % ---------------------------------------------------------
    \item \textbf{[Level C1 -- Mengingat / \textit{Remembering}]} \hfill \textbf{(Bobot: 10\%)}
    \begin{enumerate}[label=\alph*.]
        \item Sebutkan dan jelaskan secara ringkas 4 (empat) komponen dasar pembentuk deret waktu ($y_t$) dalam domain teknik elektro!
        \item Berikan masing-masing 1 (satu) contoh fenomena fisik kelistrikan yang merepresentasikan:
        \begin{itemize}
            \item Pola Tren ($T_t$)
            \item Pola Musiman Harian ($S_t$)
            \item Pola Residu/Derau Stokastik ($R_t$)
        \end{itemize}
    \end{enumerate}

    \vspace{0.3cm}

    % ---------------------------------------------------------
    % C2: MEMAHAMI (UNDERSTANDING)
    % ---------------------------------------------------------
    \item \textbf{[Level C2 -- Memahami / \textit{Understanding}]} \hfill \textbf{(Bobot: 15\%)}
    \begin{enumerate}[label=\alph*.]
        \item Jelaskan perbedaan mendasar antara model dekomposisi \textbf{Aditif} ($y_t = T_t + S_t + R_t$) dan model \textbf{Multiplikatif} ($y_t = T_t \times S_t \times R_t$) ditinjau dari perubahan variasi musiman terhadap level tren data!
        \item Mengapa kita sering menerapkan transformasi logaritmik $\tilde{y}_t = \ln(y_t)$ pada model multiplikatif sebelum melakukan pemodelan statistik linier? Tunjukkan hubungan matematisnya!
        \item Bagaimana kriteria visual untuk memilih antara model aditif vs. multiplikatif saat memeriksa grafik profil konsumsi daya listrik tahunan kawasan industri yang terus berekspansi?
    \end{enumerate}

    \newpage

    % ---------------------------------------------------------
    % C3: MENERAPKAN (APPLYING)
    % ---------------------------------------------------------
    \item \textbf{[Level C3 -- Menerapkan / \textit{Applying}]} \hfill \textbf{(Bobot: 20\%)} \\
    Tabel berikut mencatat data daya aktif (dalam MW) yang disalurkan melalui sebuah penyulang 20 kV selama 7 hari berturut-turut pada pukul 19.00 WIB:
    
    \begin{center}
    \begin{tabular}{|c|c|c|c|c|c|c|c|}
    \hline
    \rowcolor{boxbg}
    \textbf{Hari ke- ($t$)} & 1 & 2 & 3 & 4 & 5 & 6 & 7 \\ \hline
    \textbf{Beban $y_t$ (MW)} & 42.0 & 45.0 & 44.0 & 48.0 & 50.0 & 49.0 & 53.0 \\ \hline
    \end{tabular}
    \end{center}

    \begin{enumerate}[label=\alph*.]
        \item Hitunglah nilai estimasi komponen tren $\hat{T}_t$ untuk hari ke-3, 4, dan 5 ($t=3, 4, 5$) menggunakan metode \textit{Simple Centered Moving Average} dengan ordo $m = 3$ ($3$-SMA):
        \begin{equation*}
            \hat{T}_t = \frac{1}{3}(y_{t-1} + y_t + y_{t+1})
        \end{equation*}
        \item Jika diasumsikan model dekomposisi bersifat \textbf{Aditif}, hitunglah nilai komponen \textit{detrended} ($D_t = y_t - \hat{T}_t$) untuk $t=3, 4, 5$.
        \item Hitung pula nilai komponen \textit{detrended} ($D'_t = y_t / \hat{T}_t$) jika diasumsikan model \textbf{Multiplikatif}.
    \end{enumerate}

    \vspace{0.4cm}

    % ---------------------------------------------------------
    % C4: MENGANALISIS (ANALYZING)
    % ---------------------------------------------------------
    \item \textbf{[Level C4 -- Menganalisis / \textit{Analyzing}]} \hfill \textbf{(Bobot: 20\%)} \\
    Seorang insinyur proteksi dan SCADA melakukan dekomposisi aditif klasik pada data arus saluran transmisi $I(t)$ selama 1 bulan dengan periode musiman 24 jam. Setelah tren $\hat{T}_t$ dan musiman harian $\hat{S}_t$ dihilangkan, diperoleh komponen residual $\hat{R}_t$.
    
    Hasil uji statistik pada residual $\hat{R}_t$ menunjukkan:
    \begin{itemize}
        \item Nilai rata-rata $\mu_{\hat{R}} = 0.002 \approx 0$
        \item Plot fungsi autokorelasi (ACF) dari $\hat{R}_t$ masih menunjukkan koefisien autokorelasi yang signifikan pada $lag = 12$ ($r_{12} = 0.42 > \frac{1.96}{\sqrt{N}}$) dan $lag = 24$ ($r_{24} = 0.38$).
        \item Terdapat 3 lonjakan magnitudo tajam ($> 4\sigma$) yang terjadi pada tanggal 12 dan 18.
    \end{itemize}

    \textbf{Pertanyaan Analisis:}
    \begin{enumerate}[label=\alph*.]
        \item Apakah komponen residual $\hat{R}_t$ di atas sudah memenuhi kriteria sebagai proses \textbf{White Noise}? Jelaskan analisis Anda berdasarkan 3 syarat matematis \textit{white noise}!
        \item Jelaskan interpretasi fisik kelistrikan mengenai penyebab munculnya autokorelasi pada $lag = 12$ dan $lag = 24$ pada residual tersebut!
        \item Apa rekomendasi perbaikan model dekomposisi yang harus dilakukan insinyur tersebut agar informasi periodik dapat terekstraksi secara tuntas?
    \end{enumerate}

    \newpage

    % ---------------------------------------------------------
    % C5: MENGEVALUASI (EVALUATING)
    % ---------------------------------------------------------
    \item \textbf{[Level C5 -- Mengevaluasi / \textit{Evaluating}]} \hfill \textbf{(Bobot: 15\%)} \\
    Dua kelompok mahasiswa Teknik Elektro melakukan dekomposisi peramalan pada dataset daya keluaran Pembangkit Listrik Tenaga Surya (PLTS) berkapasitas 5 MWp selama 1 tahun (pencatatan interval 1 jam).
    \begin{itemize}
        \item \textbf{Kelompok A} langsung menerapkan dekomposisi Aditif klasik pada seluruh deret waktu 24 jam (termasuk malam hari ketika daya bernilai 0 kW).
        \item \textbf{Kelompok B} memisahkan data menjadi dua domain: (1) Domain Malam (18.30--05.30) diset pasti 0 kW tanpa model stokastik, dan (2) Domain Siang (06.00--18.00) dilakukan pemodelan menggunakan rasio indeks kejernihan langit (\textit{Clear-Sky Index} $k_t = P_{\text{actual}}(t) / P_{\text{clearsky}}(t)$).
    \end{itemize}

    \textbf{Tugas Evaluasi:}
    \begin{enumerate}[label=\alph*.]
        \item Evaluasilah kelemahan fatal pendekatan Kelompok A, terutama terkait estimasi komponen musiman $\hat{S}_t$ pada malam hari yang dapat menghasilkan nilai daya negatif yang tidak masuk akal secara fisik!
        \item Mengapa pendekatan \textit{Clear-Sky Index} yang diterapkan Kelompok B jauh lebih superior secara keteknikan dalam mengisolasi komponen tren dan musiman matahari?
    \end{enumerate}

    \vspace{0.4cm}

    % ---------------------------------------------------------
    % C6: MERANCANG / MEMBUAT (CREATING)
    % ---------------------------------------------------------
    \item \textbf{[Level C6 -- Merancang / \textit{Creating}]} \hfill \textbf{(Bobot: 20\%)} \\
    Untuk data deret waktu dengan frekuensi jam ($m=24$), jendela pemulus berordo genap tidak simetris pada titik tengah observasi. Oleh karena itu, standar rekayasa menggunakan filter pemulus berbobot dua-tahap: $2 \times 24$-\textit{Centered Moving Average} ($2 \times 24$-CMA) yang dirumuskan sebagai:
    \begin{equation}
        \hat{T}_t = \frac{1}{24} \left( \frac{1}{2}y_{t-12} + y_{t-11} + \dots + y_t + \dots + y_{t+11} + \frac{1}{2}y_{t+12} \right)
    \end{equation}

    \textbf{Tugas Perancangan:}
    Rancanglah sebuah \textbf{algoritma / pseudocode} lengkap (atau skrip Python/Octave) untuk melakukan dekomposisi aditif klasik dari awal (\textit{from scratch}) tanpa menggunakan fungsi bawaan pustaka \texttt{statsmodels.tsa}.
    
    Rancangan Anda harus secara eksplisit memuat:
    \begin{enumerate}[label=\alph*.]
        \item Inisialisasi array dan penghitungan $\hat{T}_t$ menggunakan formula $2 \times 24$-CMA dengan penanganan batas indeks (\textit{boundary padding / NaN handling}).
        \item Tahap \textit{detrending}: $D_t = y_t - \hat{T}_t$.
        \item Tahap kalkulasi indeks musiman rata-rata untuk setiap jam $k \in \{0, 1, \dots, 23\}$:
        \begin{equation*}
            \bar{S}_k = \frac{1}{N_k}\sum_{j} D_{24j+k}
        \end{equation*}
        dan normalisasi terpusat (\textit{centering normalization}):
        \begin{equation*}
            \hat{S}_k = \bar{S}_k - \frac{1}{24}\sum_{i=0}^{23}\bar{S}_i \quad \Longrightarrow \quad \sum_{k=0}^{23}\hat{S}_k = 0
        \end{equation*}
        \item Tahap ekstraksi residual akhir: $\hat{R}_t = y_t - \hat{T}_t - \hat{S}_t$.
    \end{enumerate}

\end{enumerate}

\vspace{0.6cm}
\begin{center}
    \small \textit{--- Selamat Mengerjakan \& Junjung Tinggi Integritas Akademik ---}
\end{center}

\end{document}
